\section*{Dane i ich przetwarzanie}
Twórcy dostarczają nam 3 datasety:
\begin{itemize}
    \item PTB-XL: \url{https://physionet.org/content/ptb-xl/1.0.3/}
    \item SaMi-Trop: \url{https://zenodo.org/records/4905618}
    \item CODE-15\%: \url{https://zenodo.org/records/4916206}
\end{itemize}

\subsection*{PTB-XL}
Badanie PTB-XL zostało przeprowadzone w Niemczech i ze względu na geograficzne położenie,
zawiera dane wyłącznie osób zdrowych. Badania EKG były oryginalnie wykonane w częstotliwości
500~Hz, ale twórcy dostarczają je również w formie resamplowanej do 100~Hz. Baza zawiera
ok.~22 tysięcy rekordów po 10 sekund z dołączonymi obszernymi metadanymi. Większość z nich
jednak nie pokrywała się z metadanymi z pozostałych dwóch badań, a np.\ przy wadze i wzroście
większość rekordów jest nieuzupełnionych. Dlatego też wstępnie uznaliśmy, że najlepiej będzie
brać pod uwagę wyłącznie wiek i płeć.

\subsection*{Sami-Trop}
Badanie zostało przeprowadzone w Brazylii i zawiera 1631 rekordów. Zawiera sygnały EKG
w częstotliwości 400~Hz po 7 lub 10 sekund. Wszyscy badani są chorzy - etykieta ta jest
pewna ze względu na przeprowadzone badania immunologiczne. Danych tabelarycznych jest
absolutne minimum: płeć, wiek, $normal\_ecg$, $death$ i $timey$.

\subsection*{CODE-15\%}
Zdecydowanie największe badanie - posiada ok.~45~GB danych. Analogicznie do SaMi-Trop,
zostało wykonane w Brazylii, dlatego sygnały EKG są w analogicznej formie. Metadanych jest
sporo, jednak zdecydowaliśmy się wyciągać z nich to, co z SaMi-Trop: płeć, wiek oraz
$normal\_ecg$. Etykiety natomiast są ,,słabe'' - o ile osoby z przypisaną jedynką
są raczej na pewno chore, to osoby uznane za zdrowe po prostu nie mają objawów choroby,
co nie świadczy o braku zakażenia ze stuprocentową pewnością.

\subsection*{Wybór baz danych}
Podczas trenowania modeli opisanych w kolejnych sekcjach napotkaliśmy na poważny problem -
modele bardzo dobrze zaczęły odróżniać dane brazylijskie od niemieckich. Mamy kilka teorii:
\begin{enumerate}
    \item Sygnały niemieckie z jakiegoś powodu miały dużo wyższą amplitudę względem tych
    brazylijskich, a standaryzacja zawiodła.
    \item Ze względu na pochodzenie etniczne, badania EKG z Niemiec bazowo różnią się od
    tych z Brazylii.
    \item Resampling zawiódł.
\end{enumerate}
Ze względu na ten kłopot, zdecydowaliśmy się całkowicie pominąć badanie PTB-XL.

\subsection{Skrypty}
Opisy użycia skryptów są na samej górze w samych plikach. Skrypty wymagają wcześniejszego ręcznego
pobrania danych na dysk (domyślnie do folderów $samitrop-data$ oraz $code15-data$).

\subsubsection*{Przygotowanie zbiorów i walidacja - skrypty $prepare\_*$}
Oba skrypty $prepare\_*$ wykonują następujące operacje na każdym recordzie:
\begin{itemize}
    \item Sygnały 7-sekundowe są dopełniane na początku i końcu zerami - dlatego pozbywamy się ich.
    \item Odrzucamy recordy zawierające NaNy i płaskie sygnały.
    \item Wycinamy losowe okna 7-sekundowe z sygnału 400 Hz.
    \item Resampling do 100 Hz.
    za pomocą $scipy.signal.resample\_poly$.
\end{itemize}

Wynikiem każdego skryptu są macierze NumPy o układzie $(N, 12, L)$, gdzie $N$ to liczba
rekordów, 12 to liczba leadów EKG, a $L$ to liczba próbek w czasie. Dodatkowe pliki zawierają etykiety, metadane
(płeć, wiek, $normal\_ecg$) oraz identyfikatory recordów.

\subsubsection*{Scalanie i przetwarzanie - skrypt $merge\_and\_process\_data.py$}
Skrypt $merge\_and\_process\_data.py$ łączy oba zbiory i wykonuje operacje:
\begin{itemize}
    \item Filtruje sygnały za pomocą filtru pasmowego 0,5 - 40 Hz. Usuwa szumy spowodowane nieistotnymi
    wydarzeniami, takimi jak oddech czy przypadkowy ruch ciała pacjenta.
    \item Aplikuje standaryzację Z - zachowuje ona kształt sygnałów dla każdego leadu,
    więc wydaje się być naturalnym wyborem.
\end{itemize}

Wynikiem końcowym są scalone macierze EKG
$ecg\_merged\_100hz*.npy$, macierz etykiet
$labels\_merged.npy$, metadanych $metadata\_merged.npy$ oraz
unikalnych identyfikatorów $ids\_merged.npy$.